% ============================================================
\subsection*{Resistencia a Ataques de Rainbow Tables}
% ============================================================

Un ataque de Rainbow Table precalcula millones de pares $\text{hash}(p_i) \rightarrow p_i$ para
un conjunto de contraseñas $\{p_i\}$. Su eficacia depende de que la misma contraseña siempre
produzca el mismo hash, condición que se cumple únicamente en ausencia de salt.

En el sistema IMC, BCrypt genera un salt aleatorio de 128~bits independiente por cada usuario.
La probabilidad de que dos usuarios con la misma contraseña $p$ generen el mismo salt es:

\begin{equation}
    P(\text{colisión de salt}) = \frac{1}{2^{128}} \approx 2{,}94 \times 10^{-39}
    \label{eq:salt_collision}
\end{equation}

Para que una Rainbow Table sea efectiva, el atacante necesitaría construir una tabla específica
para cada salt, con $2^{128}$ entradas posibles. Almacenar $2^{128}$ hashes de 24~bytes
requeriría $\approx 8{,}5 \times 10^{30}$~gigabytes --- una capacidad de almacenamiento
imposible con cualquier tecnología actual o proyectada~\cite{owasp2023}.

\begin{figure}[h]
    \centering
    \includegraphics[width=0.85\textwidth]{bcrypt_demo.png}
    \caption{Demostración en terminal: la misma contraseña \texttt{``admin123''} produce hashes
             BCrypt completamente distintos en cada llamada, evidenciando la generación de salt
             aleatorio único.}
    \label{fig:demo}
\end{figure}

\noindent En la Figura~\ref{fig:demo} se ejecutó \texttt{passwordEncoder.encode("admin123")} en dos
ocasiones consecutivas. Los hashes resultantes difieren completamente porque BCrypt deriva un nuevo
salt de 128~bits en cada invocación, confirmando la imposibilidad práctica de un ataque por
Rainbow Tables sobre el sistema IMC.
